\section{Introduction}\label{sec:intro}

\guide{Motivation: the problem, why it matters now, and who benefits.}
\filler{1}{2}

\guide{Gap: what current approaches~\parencite{placeholder-survey} cannot do, and why that is hard.}
\filler{3}{1}

\begin{figure*}[t]
  \centering
  % Placeholder teaser (from latex-tikz's pattern-pipeline.tex). Redraw it from the code with the latex-tikz skill.
% Spans the text width (figure* in the section file). Check alone: bash <latex-tikz>/scripts/tikzcheck.sh figures/overview.tex
% Embedded with \input via the standalone package, so main.tex's preamble must load everything used here.
\documentclass[tikz, border=6pt]{standalone}
\usepackage{amsmath, bm}
\usepackage{agentex-tikz}
\begin{document}
\begin{tikzpicture}[ax]
  \node[store] (data) {Input\\data};
  \node[box, right=of data] (prep) {Pre-\\processing};
  \node[box, right=of prep] (model) {Model\\$f_{\bm\theta}$};
  \node[box, right=of model] (eval) {Evaluation};
  \node[io, right=of eval] (out) {Results};
  \foreach \a/\b in {data/prep, prep/model, model/eval, eval/out}
    \draw[flow] (\a) -- (\b);
  \begin{scope}[on background layer]
    \node[group, fit=(prep)(model), label={[note]above:training stage}] {};
  \end{scope}
  \draw[skip] (eval.south) -- ++(0,-7mm) -| node[note, pos=0.25, below] {tune $\eta$} (model.south);
\end{tikzpicture}
\end{document}

  \caption{Overview of the approach. \emph{Placeholder}: replace with the paper's teaser figure, derived from the code (\texttt{latex-tikz} skill).}
  \label{fig:overview}
\end{figure*}

\guide{Approach: the key idea in two or three sentences, pointing to \cref{fig:overview}.}
\filler{4}{1}

\guide{Contributions: three or four concrete, checkable claims, each pointing to its evidence.}
\begin{itemize}
  \item A method for \dots{} (\cref{sec:method}).
  \item An evaluation on three datasets showing \dots{} (\cref{sec:results}).
  \item An analysis of when and why the method fails (\cref{sec:discussion}).
\end{itemize}

The rest of the paper is organised as follows.
\Cref{sec:related} reviews prior work, \cref{sec:method} presents the method, \cref{sec:setup,sec:results} report the experiments, and \cref{sec:discussion,sec:conclusion} discuss limitations and conclude.
